%\section*{Новембар 2025. Група А}
%\addcontentsline{toc}{section}{Нов. 2025. Група А}
\rok{Новембар 2025.}{A}

Први практични тест из Алгоритама и структура података

\zadatak{1}{Merge Sort}
%\subsubsection{Задатак 1. \textit{Merge Sort}}
%\textbf{Задатак 1. \textit{Merge Sort}} \\
Написати методу \textit{merge()} \textit{MergeSort} алгоритма за сортирање целобројног низа дужине \textit{n}.

\textbf{Додатне напомене за решавање задатка:}
\begin{itemize}
    \item Улаз је несортиран низ целих бројева.
    \item Излаз је сортирани низ целих бројева.
    \item У шаблону се налази класа \texttt{RandomArrayGenerator} коју треба искористити за генерисање низа случајних бројева.
    \item У оквиру класе \texttt{MergeSort} написати имплементацију за методу: 
    \begin{Verbatim}[fontsize=\footnotesize, numbers=left, xleftmargin=10mm]
private static void merge(long[] arr, long[] tmpArray,
                int lowerIndex, int middleIndex, int upperIndex)
    \end{Verbatim}
    
    \item Обезбедити код у \texttt{Main} методи за тестирање алгоритма.
\end{itemize}



\zadatak{2}{Спајање две једноструко спрегнуте листе, растући поредак}
Спојити две претходно сортиране, једноструко спрегнуте листе у једну, тако да елементи резултатујуће листе буду сортирани у растућем поретку. Алгоритам је неопходно написати са линеарном временском комплексношћу $O(n)$.

\textbf{Додатни захтеви за решавање задатка:}
\begin{itemize}
    \item Алгоритам писати у оквиру класе \texttt{LinkedList}.
    \item Претпоставимо да су почетне листе већ сортиране.
    \item Захтеван потпис методе: 
    \begin{Verbatim}[fontsize=\footnotesize, numbers=left, xleftmargin=10mm]
public static LinkedList MergingTwoLinkedLists(LinkedList list1, 
                                     LinkedList list2)
    \end{Verbatim}

\end{itemize}



 \textbf{Пример:}

\begin{itemize}
    \item \textbf{Улаз:} $1 \leftrightarrow 2 \leftrightarrow 3 \leftrightarrow 4$  ; \quad $1 \leftrightarrow 2 \leftrightarrow 3 \leftrightarrow 4$ \\
          \textbf{Излаз:} $1 \leftrightarrow 1 \leftrightarrow 2 \leftrightarrow 2 \leftrightarrow 3 \leftrightarrow 3 \leftrightarrow 4 \leftrightarrow 4$
    \item \textbf{Улаз:} $2 \leftrightarrow 5 \leftrightarrow 6 \leftrightarrow 7$ ; \quad $1 \leftrightarrow 4 \leftrightarrow 8 \leftrightarrow 9$ \\
          \textbf{Излаз:} $1 \leftrightarrow 2 \leftrightarrow 4 \leftrightarrow 5 \leftrightarrow 6 \leftrightarrow 7 \leftrightarrow 8 \leftrightarrow 9$
\end{itemize}

\zadatak{3}{Matryoshka Doll Problem}
Руске матријошке лутке (\textit{matryoshka}) су сет лутки које се смештају једна унутар друге. Свака лутка има одређену величину и може да садржи само лутку која је мања од ње, а може да буде смештена само у лутку која је већа од ње.

\textbf{Проблем:} Имате низ лутки различитих величина које су насумично поређане. Треба да их сортирате и угнездите тако да највећа буде споља, а унутар ње следећа по величини, итд.

\textbf{Додатни захтеви за решавање задатка:}
\begin{itemize}
    \item Задатак решавати у оквиру методе \texttt{public void MatryoshkaDollSort(int[] arr)}
    \item Затим симулирати угнеждавање користећи стекове:
    \begin{itemize}
        \item Сваки стек представља један сет угнеждених лутки.
        \item Правило: на врх стека може да стане само лутка која је мања од оне испод.
        \item Ако лутка не може да стане на тренутни стек, креира се нови стек.
    \end{itemize}
    \item Задатак решити са временском комплексношћу $O(n)$.
    \item Резултат извршавања алгоритма треба да буде сортирање оригиналног низа у опадајућем редоследу.
\end{itemize}

\textbf{Пример 1:}
\begin{itemize}
    \item \textbf{Улаз} (несортирани низ целих бројева): \texttt{[15, 6, 20, 8, 12, 4]}
    \item \textbf{Стекови:} \\
    \texttt{\{15, 6, 4\}} \quad -- 4 је на врху стека \\
    \texttt{\{20, 8\}} \quad\quad -- 8 је на врху стека \\
    \texttt{\{12\}} \quad\quad\quad -- 12 је на врху стека
    \item \textbf{Излаз 1:} \texttt{[4, 6, 8, 12, 15, 20]}
\end{itemize}

\noindent \textbf{Пример 2:}
\begin{itemize}
    \item \textbf{Улаз} (несортирани низ целих бројева): \texttt{[7, 6, 10, 5, 4]}
    \item \textbf{Стекови:} \\
    \texttt{\{7, 6, 5, 4\}} \quad -- 4 је на врху стека \\
    \texttt{\{10\}} \quad\quad\quad\quad\:\, -- 10 је на врху стека
    \item \textbf{Излаз 2:} \texttt{[4, 5, 6, 7, 10]}
\end{itemize}

\sablon{Шаблон првог теста}{2025/Decembar/GrupaA/AISP2025_Test1_GrupaA.linq}